% Generated draft: requires mathematical and source review.
\documentclass[11pt]{article}
\usepackage[margin=25mm]{geometry}
\usepackage{fontspec}
\setmainfont{Latin Modern Roman}
\usepackage{amsmath,amssymb,mathtools,unicode-math}
\setmathfont{Latin Modern Math}
\usepackage{enumitem,needspace,xurl,hyperref}
\setlength{\parindent}{0pt}
\setlength{\parskip}{5pt}
\setlength{\emergencystretch}{3em}
\newcommand{\N}{\mathbb{N}}
\newcommand{\Z}{\mathbb{Z}}
\newcommand{\Q}{\mathbb{Q}}
\newcommand{\R}{\mathbb{R}}
\newcommand{\C}{\mathbb{C}}
\newcommand{\F}{\mathbb{F}}
\newcommand{\A}{\mathbb{A}}
\newcommand{\PP}{\mathbb{P}}
\newcommand{\E}{\mathbb{E}}
\newcommand{\eps}{\varepsilon}
\newcommand{\dd}{\,\mathrm{d}}
\begin{document}
\section*{UnsolvedMath ID: 700056}
\textbf{Source problem code:} catalogue entry
\subsection{Definitions and mathematical statement}
The supplied question is: ``Is the class of groups of abelian extensions of abelian groups closed under taking direct sums $(A,B)\mapsto A\oplus B$?'' The phrase ``groups of abelian extensions'' is not defined in the supplied record, so the intended class remains unresolved.

One possible interpretation uses extension groups. For abelian groups $X$ and $Y$, let $\operatorname{Ext}^{1}_{\Z}(X,Y)$ denote the abelian group, under Baer sum, of equivalence classes of short exact sequences of abelian groups
\[
0\longrightarrow Y\longrightarrow E\longrightarrow X\longrightarrow 0.
\]
Two extensions are equivalent if there is an isomorphism between their middle groups commuting with the displayed maps and inducing the identity on $X$ and $Y$. Define the class
\[
\mathcal E=\{G:\ G\text{ is an abelian group and }G\cong\operatorname{Ext}^{1}_{\Z}(X,Y)\text{ for some abelian groups }X,Y\}.
\]
Under this interpretation, the question is whether
\[
\forall A,B\in\mathcal E,\qquad A\oplus B\in\mathcal E.
\]
Equivalently, for every four abelian groups $X_1,Y_1,X_2,Y_2$, must there exist abelian groups $X,Y$ such that
\[
\operatorname{Ext}^{1}_{\Z}(X_1,Y_1)\oplus\operatorname{Ext}^{1}_{\Z}(X_2,Y_2)
\cong\operatorname{Ext}^{1}_{\Z}(X,Y)?
\]
Here $A\oplus B$ is the direct sum of the two abelian groups.

A different interpretation concerns groups that are themselves extensions with abelian kernel and abelian quotient. Define
\[
\mathcal M=\{G:\ \text{there exists }N\trianglelefteq G\text{ such that }N\text{ and }G/N\text{ are abelian}\}.
\]
Under that interpretation the question would instead be
\[
\forall A,B\in\mathcal M,\qquad A\oplus B\in\mathcal M,
\]
where the binary external direct sum means the direct product $A\times B$ with componentwise group operation. These interpretations define different classes. The supplied wording does not establish which is intended; neither formulation is asserted here to be the definitive source problem. No finiteness or countability restrictions are supplied.
\subsection{Short English statement}
The source asks whether a certain class associated with abelian extensions is closed under binary direct sums. Its terminology leaves unresolved whether its members are Baer extension groups $\operatorname{Ext}^{1}_{\Z}(X,Y)$ or groups having an abelian normal subgroup and abelian quotient. Both conditional formulations are recorded for review.
\subsection{Sources}
\begin{itemize}
\item Local UnsolvedMath registry: supplied complete source record, permanent ID $700056$, exact source problem code \texttt{catalogue entry}, title ``Kourovka Notebook Problem 1.65.'' The registry labels the problem \texttt{open}; this is a supplied status, not an independently verified mathematical assessment.
\item Supplied attribution: L. Ya. Kulikov, proposed in 1965; \emph{The Kourovka Notebook (Unsolved Problems in Group Theory)}, first issue (1965), Problem 1.65. The record identifies the 21st issue edition (2026), arXiv:1401.0300v48. Its collection description names E. I. Khukhro and V. D. Mazurov as editors.
\item Registry provenance: upstream ID $700056$, upstream revision \texttt{d371ee917ac9647fbb06cd6fc529deec6f3dda82}; record created and updated at \texttt{2026-10-16T14:00:00Z}. These are supplied registry timestamps, not dates of independent source or literature verification.
\end{itemize}
\end{document}
